% !TeX root = ../../tesis.tex
% §2.1.4 — sustituye 2-3-Marco-Dispersion.tex l. 37-61 y 2-4-Teoria-Casos.tex
% l. 26-59 (#36). «Transiciones de Fase» (repetía 1-4) se eliminó; Aldous sale de §3.3.4.

% -----------------------------------------------------------------------
\subsection{Topología Anillo-Radial y Casos de Anillos}
\label{subsec:ea-anillo-radial}
% -----------------------------------------------------------------------

La teoría de redes ofrece dos argumentos para agregar un anillo a una red
radial. \citet{aldous2019optimal} buscan la geometría que minimiza el tiempo
promedio de viaje entre pares de puntos y encuentran que, cuando la longitud
total de la red rebasa un valor crítico $L_c$, la topología óptima cambia de
golpe y pasa a incluir un bucle o anillo. Su modelo es idealizado: indica cuándo
un anillo podría convenir, no dónde trazarlo. \citet{owais2020} llega a la misma
geometría por otra vía. Muestra que, en un sistema sobre todo radial, un anillo
reduce los transbordos necesarios para ir de una zona periférica a otra, con lo
que se acortan los traslados y el viaje multimodal pierde fricción.

\label{sec:casos-similares}
Dos metrópolis con redes radiales saturadas en el centro han optado por esa
geometría, con propósitos distintos. La Gran Línea Circular del Metro de Moscú,
inaugurada en 2023 con 70 kilómetros, se diseñó para descargar las líneas
radiales del centro y comunicar zonas periféricas sin pasar por las estaciones
centrales más congestionadas \citep{efe2023}. El \termino{Suburban Rail Loop} de
Melbourne, todavía en planeación y construcción, parte de otro diagnóstico: los
suburbios del \termino{Middle-Ring} concentran hospitales, universidades y
centros de empleo, pero no tienen conexiones laterales entre sí, y el proyecto
busca convertirlos en una red policéntrica de precintos de empleo
\citep{suburbanrailloop2021}. En Moscú el anillo resuelve un problema de
operación del centro; en Melbourne, uno de articulación de la periferia. Ninguna
de las dos fuentes mide la estructura de la red con indicadores como los de
§\ref{subsec:ea-indicadores}: una es una nota de prensa sobre la inauguración y
la otra resume el caso de negocio de un proyecto que aún no opera.

De estos casos se sabe que la geometría que esta tesis prueba sobre el Anillo
Periférico ya se eligió en otras metrópolis con problemas parecidos. No se sabe
cuánto cambia una red cuando se le agrega. Eso es lo que mide el
Capítulo~\ref{ch:analisis-red-anillo} para la \zmvm{}, con el anillo como caso
de demostración metodológica.
